\documentclass[11pt]{article}
\usepackage[T1]{fontenc}
\usepackage[english]{babel}
\usepackage[margin=1in, headheight=14pt]{geometry}
\usepackage{amsmath, amssymb, amsthm}
\usepackage{enumitem}
\usepackage{fancyhdr}

\newcommand{\course}{CODE1234 -- Course name}
\newcommand{\student}{{{AUTHOR}}}

\pagestyle{fancy}
\fancyhf{}
\fancyhead[L]{\course}
\fancyhead[R]{\student}
\fancyfoot[C]{\thepage}

\theoremstyle{definition}
\newtheorem{exercise}{Exercise}
\newenvironment{solution}{\begin{proof}[Solution]}{\end{proof}}

\begin{document}

\begin{center}
  {\LARGE\bfseries {{TITLE}} \par}
  \medskip
  \student \quad -- \quad Due \today
\end{center}

\begin{exercise}
Show that for every integer $n \ge 1$,
\[
  \sum_{k=1}^{n} k = \frac{n(n+1)}{2}.
\]
\end{exercise}

\begin{solution}
By induction on $n$. For $n = 1$, both sides equal~$1$. Assume the equality
holds for $n$. Then
\[
  \sum_{k=1}^{n+1} k = \frac{n(n+1)}{2} + (n+1) = \frac{(n+1)(n+2)}{2},
\]
which is the equality for $n+1$.
\end{solution}

\begin{exercise}
Let $f(x) = x^3 - 3x + 1$.
\begin{enumerate}[label=(\alph*)]
  \item Compute $f'(x)$.
  \item Find the local extrema of $f$.
\end{enumerate}
\end{exercise}

\begin{solution}
\begin{enumerate}[label=(\alph*)]
  \item $f'(x) = 3x^2 - 3 = 3(x - 1)(x + 1)$.
  \item $f'$ vanishes at $x = -1$ and at $x = 1$. Since $f''(x) = 6x$, $f$
        has a local maximum $f(-1) = 3$ and a local minimum $f(1) = -1$.
\end{enumerate}
\end{solution}

\end{document}
